\safechapter{Simulacro de Bloque I}
% ============================================================

Tiempo recomendado:
\[
\boxed{60\text{--}75\ \text{minutos}}
\]
sin consultar apuntes.

\Needspace{5\baselineskip}\section*{Pregunta 1 --- Teoría}
Explica brevemente:
\begin{enumerate}[label=\alph*)]
  \item diferencia entre sólido rígido y sólido deformable;
  \item homogeneidad e isotropía;
  \item principio de Saint-Venant;
  \item principio de superposición;
  \item diferencia entre esfuerzo y tensión.
\end{enumerate}

\subsection*{Respuesta orientativa}
Un sólido rígido no cambia de forma mientras que un sólido deformable sí modifica distancias relativas entre puntos.

Homogeneidad significa independencia de la posición. Isotropía significa independencia de la dirección.

Saint-Venant permite reemplazar sistemas de carga estáticamente equivalentes lejos de la zona local de aplicación.

La superposición requiere linealidad.

Un esfuerzo es una resultante sobre la sección; la tensión es una magnitud local por unidad de área.

\Needspace{5\baselineskip}\section*{Pregunta 2 --- Tensión y deformación}
Una barra de acero tiene:
\[
A=\qty{500}{\milli\metre\squared},\qquad
L=\qty{1.5}{\metre},\qquad
E=\qty{200}{\giga\pascal}.
\]
Se somete a $N=\qty{50}{\kilo\newton}$. Calcular tensión, deformación y alargamiento.

\subsection*{Solución}
\[
\sigma=\frac{50000}{500}=\boxed{\qty{100}{\mega\pascal}}
\]
\[
\varepsilon=\frac{100}{200000}=\boxed{5\times10^{-4}}
\]
\[
\Delta L=5\times10^{-4}(1500)=\boxed{\qty{0.75}{\milli\metre}}.
\]

\Needspace{5\baselineskip}\section*{Pregunta 3 --- Leyes de esfuerzos}
Una viga simplemente apoyada de $L=\qty{4}{\metre}$ soporta $q=\qty{3}{\kilo\newton\per\metre}$ uniforme.

Determina $R_A$, $R_B$, $V(x)$, $M(x)$ y $M_{\max}$.

\subsection*{Solución}
Por simetría:
\[
R_A=R_B=\frac{qL}{2}=\boxed{\qty{6}{\kilo\newton}}.
\]
\[
\boxed{V(x)=6-3x}
\]
\[
\boxed{M(x)=6x-\frac32x^2}.
\]
El máximo momento aparece cuando $V=0$:
\[
6-3x=0\Longrightarrow x=\qty{2}{\metre}.
\]
\[
M_{\max}=6(2)-\frac32(4)=\boxed{\qty{6}{\kilo\newton\metre}}.
\]

\Needspace{5\baselineskip}\section*{Pregunta 4 --- Interpretación}
Una zona de una viga tiene un diagrama de momento parabólico. ¿Qué puede afirmarse del cortante y de la carga repartida?

\subsection*{Solución}
Si $M(x)$ es de segundo grado, $V=M'$ es lineal. Además, $q=-V'$ es constante.

Por tanto:
\[
\boxed{M\text{ parabólico}\Rightarrow V\text{ lineal}\Rightarrow q\text{ uniforme}}.
\]

% ============================================================
\safechapter{Plan de estudio del Bloque I}
% ============================================================

\safesection{Día 1 --- Estructuras e hipótesis}
Domina
\[
\text{estructura real}\rightarrow\text{modelo}
\]
y las hipótesis: pequeños desplazamientos, pequeñas deformaciones, elasticidad lineal, homogeneidad, isotropía, Saint-Venant y superposición.

\safesection{Día 2 --- Tensión y deformación}
Debes poder explicar de memoria:
\[
\vect t^{(\vect n)},\qquad \sigma,\tau,
\]
\[
\varepsilon=\frac{\Delta L}{L},
\qquad
\varepsilon_{ij}=\frac12(u_{i,j}+u_{j,i}),
\qquad
\sigma=E\varepsilon.
\]

\safesection{Día 3 --- Modelo monodimensional}
Practica directriz, sección, rebanada, carga puntual, carga repartida, apoyo, articulación y empotramiento.

\safesection{Día 4 --- Esfuerzos internos}
Dibuja hasta poder hacerlo automáticamente:
\[
N^+,\qquad V^+,\qquad M^+,\qquad T^+.
\]
Practica cortes.

\safesection{Día 5 --- Leyes de esfuerzos}
